\documentclass[11pt,a4paper]{article}
\usepackage[utf8]{inputenc}
\usepackage[T1]{fontenc}
\usepackage{lmodern}
\usepackage[french]{babel}
\usepackage{amsmath,amssymb,mathtools}
\usepackage{icomma}
\usepackage{xcolor}
\usepackage{colortbl}
\usepackage{tikz}
\usepackage[most]{tcolorbox}
\usepackage{enumitem}
\usepackage{array}
\usepackage{booktabs}
\usepackage{fancyhdr}
\usepackage[hidelinks]{hyperref}
\usetikzlibrary{arrows.meta,positioning,calc,patterns,backgrounds,shapes.geometric,decorations.pathreplacing,decorations.markings}
\usepackage[a4paper,margin=2.15cm,headheight=26pt,headsep=12pt]{geometry}
\usepackage{titlesec}

% ---------------- Palette ----------------
\definecolor{primary}{RGB}{0,151,169}
\definecolor{primarydark}{RGB}{0,88,101}
\definecolor{defcol}{RGB}{0,114,178}
\definecolor{thmcol}{RGB}{0,140,120}
\definecolor{formulacol}{RGB}{198,93,9}
\definecolor{excol}{RGB}{0,135,90}
\definecolor{attncol}{RGB}{200,40,55}
\definecolor{methcol}{RGB}{90,90,100}
\definecolor{remarkcol}{RGB}{120,94,200}
\definecolor{barcol}{RGB}{0,151,169}
\definecolor{barcold}{RGB}{0,110,124}
\definecolor{modecol}{RGB}{200,55,55}
\definecolor{meancol}{RGB}{0,90,200}
\definecolor{medcol}{RGB}{160,90,190}
\definecolor{quartcol}{RGB}{0,135,90}
\definecolor{ogivecol}{RGB}{176,74,0}

% ---------------- Boîtes ----------------
\tcbset{boxbase/.style={enhanced,breakable,boxrule=0.9pt,arc=2.5pt,
  left=3mm,right=3mm,top=2.3mm,bottom=2.3mm,
  fonttitle=\bfseries,coltitle=white,toptitle=0.6mm,bottomtitle=0.6mm}}

\newtcolorbox{definition}[1][]{boxbase,colback=defcol!6,colframe=defcol,
  title={\textsc{Définition}\ifx\relax#1\relax\relax\else\textmd{ : \itshape #1}\fi}}
\newtcolorbox{propriete}[1][]{boxbase,colback=thmcol!7,colframe=thmcol,
  title={\textsc{Propriété}\ifx\relax#1\relax\relax\else\textmd{ : \itshape #1}\fi}}
\newtcolorbox{formule}[1][]{boxbase,colback=formulacol!7,colframe=formulacol,
  title={\textsc{Formule}\ifx\relax#1\relax\relax\else\textmd{ : \itshape #1}\fi}}
\newtcolorbox{exemple}[1][]{boxbase,colback=excol!6,colframe=excol,
  title={\textsc{Exemple}\ifx\relax#1\relax\relax\else\textmd{ : \itshape #1}\fi}}
\newtcolorbox{methode}[1][]{boxbase,colback=methcol!8,colframe=methcol,sharp corners,
  title={\textsc{Méthode}\ifx\relax#1\relax\relax\else\textmd{ : \itshape #1}\fi}}
\newtcolorbox{remarque}[1][]{boxbase,colback=remarkcol!7,colframe=remarkcol,
  title={\textsc{Astuce}\ifx\relax#1\relax\relax\else\textmd{ : \itshape #1}\fi}}
\newtcolorbox{attention}[1][]{boxbase,colback=attncol!7,colframe=attncol,
  title={\textsc{Attention}\ifx\relax#1\relax\relax\else\textmd{ : \itshape #1}\fi}}

% Boîte "exercice" et "corrigé" (titre libre passé en argument)
\newtcolorbox{exobox}[1]{boxbase,colback=primary!4,colframe=primary,
  before skip=8pt,after skip=8pt,title={#1}}
\newtcolorbox{corbox}[1]{boxbase,colback=excol!5,colframe=excol!85!black,
  before skip=8pt,after skip=8pt,title={#1}}

% ---------------- En-tête ----------------
\newcommand{\docpart}[1]{\def\thedocpart{#1}}
\docpart{}
\pagestyle{fancy}
\fancyhf{}
\fancyhead[L]{\small\textcolor{primarydark}{\textbf{Fiche 8} — Statistiques descriptives}}
\fancyhead[R]{\small\textcolor{primarydark}{\textsc{\thedocpart}}}
\fancyfoot[C]{\small\thepage}
\renewcommand{\headrulewidth}{0.8pt}
\renewcommand{\headrule}{\hbox to\headwidth{\color{primary}\leaders\hrule height \headrulewidth\hfill}}

\titleformat{\section}{\large\bfseries\color{primarydark}}{\thesection}{0.6em}{}
\titleformat{\subsection}{\normalsize\bfseries\color{primary}}{\thesubsection}{0.6em}{}
\setlist{topsep=2pt,itemsep=1.5pt,parsep=0pt}

% Bandeau de titre du document
\newcommand{\bandeau}[2]{%
\begin{tcolorbox}[boxbase,colback=primary,colframe=primarydark,halign=center,
   before skip=2pt,after skip=12pt]
{\LARGE\bfseries\color{white} #1}\\[3pt]{\normalsize\color{white} #2}
\end{tcolorbox}}

\newcommand{\ecm}{\sigma_m}
\newcommand{\ect}{\sigma}
\newcommand{\med}{M\!e}
\docpart{Difficile — Thème 5}
\begin{document}
\bandeau{Thème 5 — Représentations graphiques}{Niveau Difficile — exercices}

\noindent\textit{Objectif : passer du tableau au graphique complet, et raisonner sur un diagramme
circulaire comme au concours.}\medskip

\begin{exobox}{Exercice 1 — Du tableau à l'ogive et à la boîte (mise en situation)}
On mesure la \textbf{hauteur} (cm) de $N=50$ arbres d'une pépinière :
\begin{center}
\renewcommand{\arraystretch}{1.2}
\begin{tabular}{|c|c|c|c|c|c|}
\hline
Classes (cm) & $[100,120[$ & $[120,140[$ & $[140,160[$ & $[160,180[$ & $[180,200[$ \\\hline
$n_i$ & 5 & 11 & 16 & 12 & 6 \\\hline
\end{tabular}
\end{center}
\begin{enumerate}[label=\textbf{(\alph*)}]
  \item Le caractère est-il discret ou continu ? Quel graphique choisir, et quelle est la classe modale ?
  \item Construire les couples $(\text{borne supérieure},\ f_c)$ nécessaires pour tracer l'ogive.
  \item Déterminer la médiane par interpolation linéaire. Entre quelles graduations la lirait-on
        sur l'ogive ?
  \item On donne $Q_1\approx133{,}6$ et $Q_3\approx169{,}2$. Donner les cinq nombres de la boîte à
        moustaches ($x_{\min},Q_1,\med,Q_3,x_{\max}$) et calculer l'écart interquartile.
\end{enumerate}
\end{exobox}

\begin{exobox}{Exercice 2 — Raisonner sur un diagramme circulaire}
Dans une faculté de médecine de $3000$ étudiants, il y a trois sections. La section \textbf{Médecine}
occupe exactement la \emph{moitié} du disque ; le secteur \textbf{Sciences biomédicales} est plus
petit que le secteur \textbf{Dentisterie}.
\begin{center}
\begin{tikzpicture}[scale=1]
  \def\r{1.7}
  \fill[yellow!70,draw=black,line width=0.5pt] (0,0) -- (90:\r) arc (90:270:\r) -- cycle;
  \fill[black!75,draw=black,line width=0.5pt] (0,0) -- (90:\r) arc (90:18:\r) -- cycle;
  \fill[primary!70,draw=black,line width=0.5pt] (0,0) -- (18:\r) arc (18:-90:\r) -- cycle;
  \node[font=\scriptsize] at (-0.7,0){Médecine};
  \node[font=\scriptsize,white] at (0.85,0.75){Sc.\ bio.};
  \node[font=\scriptsize,white] at (0.75,-0.7){Dentist.};
\end{tikzpicture}
\end{center}
\begin{enumerate}[label=\textbf{(\alph*)}]
  \item Combien d'étudiants la section Médecine compte-t-elle ?
  \item Que vaut la somme des effectifs des deux autres sections ?
  \item Parmi les valeurs $600$, $900$, $1500$, $2100$, laquelle \emph{peut} correspondre au nombre
        d'étudiants \emph{en plus} en Médecine qu'en Sciences biomédicales ? Justifier en éliminant
        les autres.
\end{enumerate}
\end{exobox}
\end{document}
